(1) 根据平行轴定理，悬挂支点的转动惯量 $I_O=I_G+mh^2$。取偏离竖直向下的角度为 $\theta$，重力矩为 $-mgh\sin\theta$，刚体角动量方程给 $(I_G+mh^2)\ddot\theta=-mgh\sin\theta$；小角度时用 $\sin\theta\simeq\theta$。
(2) 线性简谐运动角频率 $\omega^2=mgh/(I_G+mh^2)$，所以周期 $T=2\pi\sqrt{(I_G+mh^2)/(mgh)}$。与长度 $\ell$ 的单摆公式 $T=2\pi\sqrt{\ell/g}$ 比较，得等效长度 $\ell_{\rm eq}=I_G/(mh)+h$。
(3) 设 $k^2=I_G/m$，则 $\ell_{\rm eq}=h+k^2/h$。将支点换至质心另一侧、距离 $h'=k^2/h$ 处，新等效长度 $\ell_{\rm eq}'=h'+k^2/h'=k^2/h+h=\ell_{\rm eq}$，故两支点的周期一致；两支点间距 $h+h'$ 恰为等效摆长，称复摆可逆性。
(4) 因为 $h+k^2/h\ge2k$（算术几何平均不等式），当 $h=k=\sqrt{I_G/m}$ 时取最小等效长度 $\ell_{\rm eq,min}=2\sqrt{I_G/m}$。此时最小周期为 $2\pi\sqrt{2\sqrt{I_G/m}/g}$；极小与极大 $h$ 均会使周期增长，符合公式。
